\cdcsection{Délais et livrables}

\cdcsubsection{Planning de réalisation}

Le projet se déroule en six phases successives. Chaque phase se termine par un jalon validé avec l’encadrant avant le passage à la suivante. Les dates des jalons sont celles du calendrier du module.

\begin{cdctableau}{P{0.9cm}P{3.6cm}LP{3.4cm}}{\cdcth{Ph.} & \cdcth{Phase} & \cdcth{Contenu} & \cdcth{Jalon de fin}}
  P1 & Cadrage & Analyse du besoin, rédaction et validation du cahier des charges & Cahier des charges validé \\
  P2 & Conception & Modélisation UML, architecture, schéma de la base de données, maquettes détaillées & Revue de conception \\
  P3 & Développement de la version~1 & Fonctions essentielles (section~\ref{sec:perimetre}) en itérations courtes & Démonstration de la version~1 \\
  P4 & Recette de la version~1 & Exécution du cahier de recette, correction des anomalies & Version~1 validée \\
  P5 & Développement de la version~2 & Fonctions avancées & Démonstration de la version~2 \\
  P6 & Recette finale et clôture & Recette de la version~2, documentation, préparation de la soutenance & Soutenance du projet \\
\end{cdctableau}

\cdcsubsection{Livrables intermédiaires}

Chaque phase produit des livrables intermédiaires, soumis à l’encadrant au jalon correspondant.

\begin{cdcitems}
  \item \textbf{Cahier des charges} (P1). Le présent document.
  \item \textbf{Dossier de conception} (P2). Diagrammes de cas d’utilisation, de classes et de séquence, architecture et modèle de données.
  \item \textbf{Maquettes détaillées} (P2). Écrans principaux de chaque profil d’utilisateur.
  \item \textbf{Version~1 de la plateforme} (P3). Code source, base de données et plateforme déployée avec les fonctions essentielles.
  \item \textbf{Cahier et procès-verbal de recette} (P4, P6). Cas de test liés aux exigences et résultats obtenus.
  \item \textbf{Version~2 de la plateforme} (P5). Plateforme enrichie des fonctions avancées.
  \item \textbf{Documentation} (P6). Manuel utilisateur par profil et guide d’installation.
  \item \textbf{Rapport final et soutenance} (P6). Rapport du projet et support de présentation.
\end{cdcitems}

Les procès-verbaux de réunion sont produits tout au long du projet (annexe~A).
